\section{Problem Formulation}
\label{sec:formulation}
\subsection{Fixed-Horizon Rest-to-Rest Task}
Consider $N$ labeled UAVs in a known workspace $\mathcal W=[\underline{\mathbf w},\overline{\mathbf w}]\subset\R^3$. Static obstacles are closed axis-aligned bounding boxes (AABBs), $\mathcal O_b=[\underline{\mathbf o}_b,\overline{\mathbf o}_b]$, $b=1,\ldots,B$. Each UAV is represented by a sphere of radius $r$, with a required surface clearance $c$. Task identities are fixed. UAV $i$ moves from its assigned start $\mathbf s_i$ to goal $\mathbf g_i$ over the prescribed common horizon $T_f$, with
\begin{equation}
 \bp_i(0)=\mathbf s_i,\quad \bp_i(T_f)=\mathbf g_i,\quad
 \bv_i(0)=\bv_i(T_f)=\mathbf0 .
 \label{eq:endpoint}
\end{equation}
Both endpoints are inside the free workspace and mutually separated. In the following, $\bv_i^0=\bv_i^f=\mathbf0$ denote these prescribed endpoint velocities. Vector inequalities are understood componentwise.

\subsection{Motion Model and Reduced Variables}
The planning-level translational dynamics are
\begin{equation}
 \dot{\bp}_i=\bv_i,\qquad \dot{\bv}_i=\ba_i.
 \label{eq:continuous}
\end{equation}
Let $h=T_f/K$ and $t_k=kh$, $k=0,\ldots,K$, and take acceleration to be constant on each interval. Exact integration yields
\begin{subequations}\label{eq:discrete}\begin{align}
 \bv_{i,k+1}&=\bv_{i,k}+h\ba_{i,k},\label{eq:velocityupdate}\\
 \bp_{i,k+1}&=\bp_{i,k}+\frac h2(\bv_{i,k}+\bv_{i,k+1}).\label{eq:positionupdate}
\end{align}\end{subequations}
Eliminating acceleration reduces the physical decision fields to $\mathbf P,\mathbf V$:
\begin{equation}
 \ba_{i,k}=\frac{\bv_{i,k+1}-\bv_{i,k}}h,\qquad
 \mathbf E_{i,k}=\bp_{i,k+1}-\bp_{i,k}
 -\frac h2(\bv_{i,k+1}+\bv_{i,k}).
 \label{eq:defect}
\end{equation}
An exact discrete trajectory has $\mathbf E_{i,k}=\mathbf0$. During optimization, this relation is relaxed to allow the position and velocity fields to adjust jointly. Section~\ref{sec:arcs} gives the final control-integration procedure.

The Euclidean speed and acceleration limits are $v_{\max}$ and $a_{\max}$. Cartesian limits $v_c=v_{\max}/\sqrt3$ and $a_c=a_{\max}/\sqrt3$ give
\begin{equation}
 -v_c\one\leq\bv_{i,k}\leq v_c\one,\qquad
 -a_c\one\leq\ba_{i,k}\leq a_c\one .
 \label{eq:bounds}
\end{equation}
These componentwise bounds are conservative sufficient conditions for the Euclidean limits and preserve linear constraints. This model describes translational motion with continuous velocity and piecewise-constant acceleration.

\subsection{Workspace and Collision Conditions}
For $\bp\in\R^3$, define the component gap to an obstacle box as
\begin{equation}
 \mathbf d_b(\bp)=\max\{\underline{\mathbf o}_b-\bp,\,
 \bp-\overline{\mathbf o}_b,\,\mathbf0\}.
 \label{eq:gap}
\end{equation}
The physical requirements, imposed over the complete horizon, are
\begin{subequations}\label{eq:avoidance}\begin{align}
 \underline{\mathbf w}+r\one&\leq\bp_i(t)\leq\overline{\mathbf w}-r\one,\\
 \norm{\mathbf d_b(\bp_i(t))}_2&\geq r+c,\qquad \forall b,\\
 \norm{\bp_i(t)-\bp_j(t)}_2&\geq 2r+c,\qquad \forall i<j.
\end{align}\end{subequations}
The clearance conditions apply to the entire motion of the spherical planning footprint.

\subsection{Acceleration-Effort Objective}
With task completion fixed by~\eqref{eq:endpoint}, the nominal objective is
\begin{equation}
 C(\mathbf V)=w_a h\sum_{i=1}^N\sum_{k=0}^{K-1}
 \left\|\frac{\bv_{i,k+1}-\bv_{i,k}}h\right\|_2^2,\qquad w_a>0.
 \label{eq:cost}
\end{equation}
This control-effort surrogate discourages large accelerations during local refinement. The flight duration $T_f$ is prescribed independently of computation time. The outer procedure terminates when recovery and trajectory acceptance are achieved within the stopping rule; the objective regularizes the recovered motion. Global optimality of the original nonconvex problem is not assumed.

The nominal problem minimizes~\eqref{eq:cost} subject to exact dynamics, fixed boundary states, motion limits, and continuous collision avoidance. Geometric avoidance makes this problem nonconvex. Spatial and relative corridors provide conservative local linear restrictions, and the recovery mechanism handles their compatibility with the motion model.
